% PDF-derived LaTeX transcription. Original manuscript source was unavailable.
% Published PDF SHA-256: 687e1c0c5a222b27674ccdd249dca3de88bf446a262946dc9ee70085e195e3f0
% Extracted text SHA-256: d6ba03c3218c76b81c439bdf46cb09b3fc66022f8058450b2d4ed636b58ba270
\documentclass[11pt]{article}
\usepackage[margin=1in]{geometry}
\usepackage{amsmath,amssymb}
\setlength{\emergencystretch}{3em}
\title{Peak and Spectral Agreement Do Not Establish Field Reconstruction}
\author{}
\date{}
\begin{document}
\maketitle
\noindent\textbf{PDF-derived LaTeX transcription.} This editable text follows the PDF extraction. Refer to the published PDF for figures, table layouts, and mathematical notation.

% Extracted chunk 1

Peak and Spectral Agreement Do Not Establish Field
Reconstruction

\section*{Abstract}

Scientific-field codecs can appear accurate under a selected diagnostic while failing a dif-
ferent reconstruction contract. We examine saved measurements from a selected preliminary
hypersonic peak-fitting test, a distinct linked full-training quality assessment, and reactor-field
spectral variants. In the preliminary test, the predicted maximum kinetic-energy density was
1.06938784 times the target maximum, but mean kinetic-energy-density retention was only
0.04041034; field and divergence diagnostics remained adverse. The linked full-training check-
point separately failed its scalar-speed spectral gate. In reactor-field panels, changing from a
multiframe vector-spectrum objective to an audit-aligned final-frame speed objective was associ-
ated with improved fitted slopes: a stronger aligned variant reached a mean slope of \ensuremath{-}3.4880426
against \ensuremath{-}3.4090753 for the reference. Its continuity-residual discrepancy and energy error nev-
ertheless remained above their recorded gates. Wavelet-tail variants approached the reference
slope more closely while worsening these other diagnostics. Equations and finite mathematical
examples explain the non-equivalence of peak, mean, power, slope, directional-field and con-
tinuity measurements. These retrospective results support explicit metric contracts and joint
reporting, not independent generalization or physically validated reconstruction.

\section*{1 Introduction and contributions}

A scientific-field codec reconstructs more than a scalar summary. Its output may be evaluated
through localized extrema, bulk energy, spatial alignment, vector direction, spectral distribution
and discrete balance residuals. These measurements answer different questions. A peak can be
recovered while the surrounding field is nearly absent; a fitted spectral exponent can agree while
amplitude or direction is wrong. Consequently, interpreting a successful diagnostic requires speci-
fying its observable, support, reduction and reference population.
This article presents a bounded empirical case study of these distinctions. The evidence comprises
saved numerical reports and captured source definitions, rather than a new architecture or new
physical experiment. The central comparisons are a selected preliminary hypersonic peak-fitting
test (hereafter, the smoke), a distinct full-training checkpoint assessment, and reactor-field variants
whose spectral training objective changed toward the evaluator’s definition. We call the latter
population CSTR, following the recorded thermal-runaway benchmark, and retain A123 as the
companion slurry benchmark name.
The contribution is to connect the reported measurements to their actual mathematical contracts.
We distinguish maximum kinetic-energy density from mean kinetic-energy density and total energy,
vector spectra from scalar-speed spectra, and global mass summaries from normalized continuity-
residual discrepancies. We also examine the limits of the nominal evaluation population after
adaptation. This study shares tokenizer lineage with a separate periodic fluid/magnetic harmonic-
decoder study, but its principal measurements concern hypersonic peak-versus-mean reconstruction
1

and CSTR spectral-objective alignment, not harmonic recovery or discrete curl constraints. The
scientific outcome is not a qualified codec: apparently favorable peak and slope measurements
coexist with important reconstruction failures.

\section*{2 Related work}

WIRE uses complex Gabor activations for visual implicit neural representations [ 1]. Its inspected
formulation preserves complex-valued intermediate representations and takes the real part of the
output for real signals. This provides relevant architectural context, but the captured implemen-
tation here is not identical to complex WIRE. Despite its class name, its activation is the real
Gaussian-windowed cosine
a(x) = exp
[U+0000]
\ensuremath{-}(\ensuremath{\sigma}x)2[U+0001]
cos(\ensuremath{\omega}x).
We therefore describe this decoder as Gabor-inspired. The captured decoder also includes coordi-
nate Fourier features, latent/metadata conditioning and dyadic residual heads. The cited visual
reconstruction results do not establish shock reconstruction or conservation for this implementation.
Focal Frequency Loss compares complex Fourier coefficients with adaptive frequency weighting
in image reconstruction and synthesis [ 2]. Its coefficient distance retains amplitude and phase
information; the dynamic weights emphasize currently difficult frequencies. This is conceptually
different from reducing power into radial shells and fitting a scalar slope. Focal Frequency Loss is
not an empirical comparator evaluated in the present evidence, and its image results do not validate
physical shocks. Neither reference supplies the benchmark measurements reported below. Our use
of these references is restricted to the inspected full-text formulations.

\section*{3 Scientific methods and validation protocols}

Evidence and populations
This is a retrospective reanalysis of saved summaries, captured computational definitions and cur-
rent manifest metadata. No additional neural training, neural inference or physical simulation was
performed. Later local arithmetic and finite-array tests are distinguished from the original codec
experiments throughout.
The peak-fitting smoke used one selected hypersonic sample and 5,000 optimization steps. The
linked full-training record instead reports 50,000 steps on 72 ranks, with 78,000 fluid rows and 70,000
graph rows. These 148,000 mixed rows are not 148,000 fluid trajectories. Its quality assessment
selects training-manifest examples; it does not constitute full-population qualification. Neither its
checkpoint nor its selected evaluation population is identified with the smoke.
The CSTR candidate reports use 323 grids, 16 frames and 8 channels, with a 2,048-value learned
latent and a 65-value metadata header. The evaluator supplies or rebuilds metadata from the target
reference and conditions physical decoding on it. These measurements are therefore not latent-only
reconstruction, and the dimensions do not establish a bit-compression ratio. Saved configurations
use bfloat16 training; the captured evaluation path encodes under bfloat16 autocasting on CUDA
and casts latents, decoded fields and metric inputs to float32. The earlier vector-spectrum helper
disables autocasting and uses float32 unless given float64 input.
The inspected source generator prescribes analytic synthetic A123 slurry and Aarti CSTR fields
over normalized spatial and temporal coordinates, with jamming, runaway and cooling envelopes.
2

These are formula-defined benchmark tensors, distinct from a transported-state advection–diffusion–
reaction generator described in a companion reactor study. The captured definition does not au-
thenticate execution of that exact source snapshot in every historical quality-assessment run. Exact
historical input arrays and source-execution bindings remain unavailable; no resolved reactor dy-
namics are inferred from the CSTR labels. Source and companion-version bindings are retained
with Supplementary record S1.
An earlier captured header has six interleaved statistics per channel—mean, standard deviation,
minimum, maximum and lower/upper quantiles—plus five aggregate values. With eight channels
this gives 53 values. Its affine normalization uses the per-channel mean and standard deviation,
with a 10\ensuremath{-}6 standard-deviation floor. This source establishes a header convention, not the identity
of every historical decoder diagnostic.
Peak, mean and field measurements
Let vc(t, x) denote the velocity components and r(t, x) the raw density channel. Throughout, hats
denote reconstructed quantities and unhatted quantities denote their references. The captured
kinetic-energy-density convention is
\ensuremath{\rho} = softplus(r), e (t, x) = 1
2 \ensuremath{\rho}(t, x)
3X
c=1
vc(t, x)2.
For the space-time support Ω of a case, define
emax = max
i\ensuremath{\in}Ω
ei, ¯e = 1
|Ω|
X
i\ensuremath{\in}Ω
ei, E \ensuremath{\Sigma} =
X
i\ensuremath{\in}Ω
ei.
These are respectively a maximum, mean and discrete total. The total is a sample sum, not an
independently calibrated physical integral. Equal prediction/target sample counts make their mean
and sum ratios equal, but neither ratio follows from a maximum ratio.
For fraction f , let Sf contain the kf = max(1 , \ensuremath{\lfloor}f |Ω|\ensuremath{\rfloor}) highest-energy target cells, selected over
flattened space and time. The peak relative-L1 objective is
L1,f = 1
kf
X
i\ensuremath{\in}Sf
|bei \ensuremath{-} ei|
max(|ei|, 1) .
The reported narrow support uses f = 0.01. The wider objective instead uses f = 0.05 and
Llog,f = 1
kf
X
i\ensuremath{\in}Sf
[log(1 + max(bei, 0)) \ensuremath{-} log(1 + max(ei, 0))]2 .
Thus the narrow relative loss and wider logarithmic loss have different supports and penalties. A
gate correlation describes agreement with a gate target, not field accuracy. We retain the smoke
report’s field relative-L2 and unmasked divergence MSE as separate reported diagnostics; their raw
values have no independently established SI calibration here.
Spectral contracts
For a field g on a cubic grid with N cells per axis, write Fg for a full spatial fast Fourier transform
(FFT); // denotes integer division. Given a shell assignment b(k), define the shell mean
Pj =
P
b(k)=j Q(k)
max(nj, 1) , n j = \#\{k : b(k) = j\}.
3

The old objective operates on multiframe vector fields. Each spatial transform uses all velocity
channels, with
Qold(k) = 1
2
3X
c=1
|Forthovc(k)|2.
Time contributes multiple fields to the training reduction rather than selecting only the final frame.
With B = max(4, N//2) and radius rk from FFT frequencies, its shells are
bold(k) = clamp
[U+0012][U+0016] rk
max(rmax, 10\ensuremath{-}12) (B \ensuremath{-} 1)
[U+0017]
, 0, B \ensuremath{-} 1
[U+0013]
.
The fitted window begins at a = max(1 , min(2, B \ensuremath{-} 2)) and ends, exclusively, at b = max( a +
2, min(B, \ensuremath{\lceil}0.85B\ensuremath{\rceil})).
The corrected observable is final-frame scalar speed,
s(x) =
3X
c=1
vc(T, x)2
!1/2
, Q new = |F(s \ensuremath{-} ¯s)|2,
using an unnormalized FFT. For integer-frequency coordinates q = N fftfreq(N ) on each axis,
bnew(q) = clamp(round \ensuremath{\parallel}q\ensuremath{\parallel}, 0, N//2).
The audit fits shells j \ensuremath{\in} [\ensuremath{\ell}, h), where
\ensuremath{\ell} = max(2, N//10), h = max(\ensuremath{\ell} + 4, \ensuremath{\lfloor}0.45N \ensuremath{\rfloor}).
Writing xj = log max(j, 1) and yj = log max(Pj, 10\ensuremath{-}12), the fitted exponent is
\ensuremath{\beta} =
P
j(xj \ensuremath{-} ¯x)(yj \ensuremath{-} ¯y)
max\{P
j(xj \ensuremath{-} ¯x)2, \ensuremath{\varepsilon}x\} .
The audit uses \ensuremath{\varepsilon}x = 10 \ensuremath{-}12. The newer training loss uses 10\ensuremath{-}8, bounds the upper slice by the
available shell count, and reduces squared prediction/target slope differences across cases. It follows
the audit definition but is not byte-identical to it.
Table 1: Spectral contracts for the CSTR comparison. All spectra are shell means, not shell-
integrated physical energy.
Property Old training objective Audit-aligned objective
Observable Multiframe velocity vector Final-frame scalar speed
FFT Full, orthonormal Full, unnormalized
Mean removal Not in captured helper Spatial speed mean removed
Power Half channel-summed squared
magnitude
Squared scalar magnitude
Shells Radius-normalized, floored Integer radius,
rounded/clamped
Fit window [a, b) defined above [\ensuremath{\ell}, h) defined above
Power floor 10\ensuremath{-}12 10\ensuremath{-}12
4

The linked full-training QA has yet another implementation; in its fit window below, \ensuremath{\ell} denotes its
own lower endpoint, max(2, \ensuremath{\lfloor}0.12B\ensuremath{\rfloor}), rather than the CSTR lower endpoint. On final-frame scalar
speed at 1283, it removes the mean, uses an unnormalized FFT and radius-normalized floored shell
means with B = N//2. It fits by least squares with an intercept over [max(2, \ensuremath{\lfloor}0.12B\ensuremath{\rfloor}), max(\ensuremath{\ell} +
4, \ensuremath{\lfloor}0.55B\ensuremath{\rfloor})), using a 10\ensuremath{-}20 power floor. It must not be substituted for the CSTR audit definition.
Continuity, energy and aggregation
For the CSTR audit, each raw density value is retained if nonnegative and softplus-transformed if
negative; the resulting density is then floored at 10\ensuremath{-}6. This elementwise rule differs from applying
softplus everywhere in the smoke energy convention. Let DMAC denote the discrete divergence of
face fluxes on a staggered marker-and-cell (MAC) grid. Define the cellwise discrete residual
Rn = (\ensuremath{\rho}n+1 \ensuremath{-} \ensuremath{\rho}n) + DMAC(\ensuremath{\rho}up
n un+1/2).
The implementation averages adjacent cell velocities onto interior faces, midpoints velocity in time,
upwinds density from the earlier time, and uses an open positive- x outlet with ambient inflow
density 1; other exterior velocity faces are zero. Its time-indexed normalized residual is
qn = \ensuremath{\langle}|Rn|\ensuremath{\rangle}x
max(\ensuremath{\langle}|\ensuremath{\rho}n+1 \ensuremath{-} \ensuremath{\rho}n|\ensuremath{\rangle}x, 1) + max(\ensuremath{\langle}|DMAC(\ensuremath{\rho}up
n un+1/2)|\ensuremath{\rangle}x, 1) .
The quantity recorded as mass\_loss\_pct is
C = 100 \ensuremath{\langle}|bqn \ensuremath{-} qn|\ensuremath{\rangle}n
max(\ensuremath{\langle}|qn|\ensuremath{\rangle}n, 1) .
We call it normalized continuity-residual discrepancy, not the fraction of total mass lost. Default
timestep and spacing are effectively 1 in this difference contract; these constants and field semantics
lack independent physical calibration.
The CSTR energy diagnostic adds a density-weighted nonnegative temperature/enthalpy proxy
to the kinetic sample sum. When available, the evaluator substitutes the decoded enthalpy head
into that field before computing energy. Its error is
H = 100 |bE \ensuremath{-} E|
max(|E|, 10\ensuremath{-}8) .
This domain-specific diagnostic is not the full-training QA’s global kinetic-energy comparison.
Wavelet-tail supervision uses relative Haar detail-band residual MSE, normalized by target detail
power floored at 10\ensuremath{-}8, averaged across bands and two levels, and then averaged between final-frame
velocity and speed losses.
For case i, spectral relative error is
di = |b\ensuremath{\beta}i \ensuremath{-} \ensuremath{\beta}i|
max(|\ensuremath{\beta}i|, 10\ensuremath{-}6) .
Saved summaries compute casewise means and float32 95th percentiles. For the candidate panels,
the dimensionless temporal-spike diagnostic is the maximum squared velocity second difference
divided by the larger of its mean and 10\ensuremath{-}12. The second difference is taken across consecutive
frames, its squared components are summed, and the maximum and mean are taken over space
5

and time within each case. Unlike the full-training temporal proxy, this diagnostic does not spatially
average velocity before differencing. Recorded percentile gates are C \ensuremath{\leq} 2\%, H \ensuremath{\leq} 10\% and d \ensuremath{\leq} 0.05;
these are evaluation contracts, not scientist-validated physical tolerances. A mean of casewise ratios
is not a ratio of mean slopes. We use saved di aggregates rather than reconstructing them from
aggregate slopes.
Finite mathematical validation
Later local tests used float64 arrays, ten random cases across 163 and 243 grids, and four analytic
examples, with a predeclared 10\ensuremath{-}10 arithmetic tolerance. They tested gain versus slope and energy,
sign and translation invariance, peak versus mean, and vector versus speed observables under the
stated discrete conventions. These tests, implemented with NumPy 2.3.5, check mathematical
non-implications, not original model inference.
For fixed density and nonzero gain a, bv = av multiplies kinetic energy and speed power by
a2. With inactive power floors, logbPj = log Pj + 2 log |a|, so the fitted slope is unchanged. A sign
reversal also preserves speed while changing vector direction. A periodic translation changes Fourier
phase but not power, and can create nonzero aligned field error. Finally, v(x) = (cos x, sin x, 0)
has constant speed but nonconstant vector components: its mean-removed speed spectrum vanishes
although its vector spectrum does not. These standard finite examples explain why a fitted exponent
is neither field recovery nor inertial-range-law validation.

\section*{4 Experimental results}

Peak fitting without bulk reconstruction
The selected hypersonic smoke achieved close agreement on its maximum and narrow target-selected
tail, but not on the bulk field. Table 2 reports the post-smoke diagnostics together to avoid
interpreting peak agreement as general recovery.
The maximum ratio therefore describes a successful extremum fit in this selected case, not ap-
proximate recovery of total energy. The reconstructed mean kinetic-energy density was only a small
fraction of the reference mean, and the reported field relative-L2 remained near unity. The maxi-
mum divergence overshoot and unmasked divergence MSE provide additional adverse diagnostics
under the recorded convention. Velocity divergence need not vanish in compressible flow, however,
and these measurements are not independently calibrated tests of the density-weighted continuity
equation or physical conservation. The well-correlated gate and small narrow-tail loss therefore do
not override the poor bulk and field diagnostics.
Distinct full-training QA failure
The full-training summary reports completion PASS. Its exact linked checkpoint QA reports F AIL:
decoded and reference final-frame scalar-speed slopes were \ensuremath{-}4.1108651 and \ensuremath{-}6.3195229, respec-
tively, producing spectral relative error 0.34949755 against 0.05. This measurement used 1283
fields and the distinct full-training spectral implementation described above.
The same QA reported passing global mass and kinetic-energy summaries. These global-sum
checks do not contradict spectral failure, nor do they satisfy the later CSTR continuity-residual
contract. The rotation consistency check was explicitly an invariance proxy, and the text-fragment
6

Table 2: Post-smoke diagnostics for the selected hypersonic case after 5,000 steps. Energy-density
and divergence values are in the captured numerical convention, without independently established
SI calibration.
Diagnostic Value
Predicted maximum kinetic-energy density 39,750,408
Target maximum kinetic-energy density 37,171,180
Maximum ratio 1.06938784
Predicted mean kinetic-energy density 544,259.5625
Target mean kinetic-energy density 13,468,323
Mean retention ratio 0.04041034
Mean relative error 0.959589656
Gate correlation 0.996647
Top-1\% relative L1 0.0231164
Top-5\% log-energy MSE 14.80047
Reported field relative L2 0.99999976
Maximum absolute divergence ratio 10.4145
Unmasked divergence MSE 4146.985
check was a proxy rather than atomistic topology validation. Completion and those proxy passes
do not establish full scientific qualification.
CSTR spectral alignment and remaining failures
Table 3 compares the seven candidate panels central to the old-versus-aligned comparison. Names
describe the recorded variants rather than implying a controlled intervention. The old-vector
lower/higher-weight and audit-aligned moderate/strong variants each record 1,500 continuation
steps, with slope-loss weights 10/30 and 30/100, respectively; all four record QA mass and energy
weights of 100 and 30. The three wavelet variants each record 2,000 continuation steps, slope-
loss weights of 40, 60 and 35 in table order, and zero QA mass and energy weights. These are
configuration-specific loss multipliers, not comparable effect sizes across the different objectives.
Continuation ancestry and other settings also differ; variant crosswalks are retained in Supplemen-
tary record S1. The CSTR reference mean slope was \ensuremath{-}3.4090753. All entries are saved casewise
aggregates; rounding in the table does not alter gate assignments.
The old vector variants remained near \ensuremath{-}5.26 under the final-frame speed audit despite their slope-
training objective. The strong aligned variant reached \ensuremath{-}3.4880426 with mean spectral relative error
0.02316399. Its CSTR spectral 95th percentile passed the recorded gate. Continuity discrepancy
and energy error did not: their means were 15.32657\% and 18.11088\%, and their 95th percentiles
were 15.67740\% and 18.18186\%.
The mild- and stronger-coupling wavelet variants reached slopes near \ensuremath{-}3.412, yet had larger
continuity and energy discrepancies than the strong aligned variant. Their spectral agreement is
therefore a tradeoff in these records, not evidence of jointly accurate thermodynamic reconstruction.
The mixed-manifest wavelet variant also passed the CSTR spectral gate while failing the other
contracts.
All seven full panels reported overall F AIL and failure count 200. This is not a statement
that every CSTR spectral check failed. The companion A123 branch remained spectrally outside
7

Table 3: CSTR results within reused candidate panels. Each panel contains 100 CSTR and 100
A123 rows. ¯\ensuremath{\beta} is mean decoded slope, ¯d and d95 summarize casewise spectral relative error, and
¯C, ¯H are mean continuity discrepancy and energy error in percent.
Variant ¯\ensuremath{\beta} ¯d d95 ¯C ¯H
Old vector, lower weight \ensuremath{-}5.261656 .543427 .545194 14 .81017 11 .46261
Old vector, higher weight \ensuremath{-}5.265267 .544486 .546238 15 .29098 12 .11382
Audit-aligned, moderate \ensuremath{-}3.852128 .129963 .132730 16 .25037 18 .60241
Audit-aligned, strong \ensuremath{-}3.488043 .023164 .026334 15 .32657 18 .11088
Wavelet, mild coupling \ensuremath{-}3.412379 .001199 .002866 27 .72021 38 .56628
Wavelet, mixed manifest \ensuremath{-}3.420659 .003398 .005109 29 .98826 34 .89359
Wavelet, stronger coupling \ensuremath{-}3.412995 .001413 .003414 25 .40436 37 .61141
the gate, with spectral 95th percentile approximately 0.182923; its energy 95th percentile was
11.35641\%, and its temporal-spike 95th percentile was 232.88837 against a recorded threshold of
100. The saved configurations froze that branch, and the panels reuse the nominal benchmark
rather than constituting independent replications.
Current manifest overlap and reference consistency
The current metadata audit inspected 3,000 manifest rows: a 1,000-row fine-tuning manifest and
a 2,000-row manifest labeled out-of-distribution (OOD). The fine-tuning manifest exactly matches
the 1,000 CSTR records in the latter. Every reference path among the first 100 CSTR evaluation
rows appears in that current training manifest. The first 100 A123 reference paths do not appear
in the CSTR-only file, but ancestral-training independence is unknown.
The current first-100 reference-slope means, accumulated in float64, are \ensuremath{-}3.4090752959 for CSTR
and \ensuremath{-}6.1556583834 for A123. They agree with the saved float32 aggregate means \ensuremath{-}3.4090752602
and \ensuremath{-}6.1556587219 within the predeclared absolute tolerance 10\ensuremath{-}6 max(1, |recorded mean |). This
supports consistency of current metadata with the saved slope summaries, not exact historical-array
identity. The current OOD manifest modification time postdates the historical candidate audits;
the original historical manifest hashes and arrays have not been authenticated.

\section*{5 Discussion and limitations}

The observations illustrate a hierarchy of information loss rather than a single universal ranking of
metrics. A maximum constrains an extremum but leaves its occupied volume unspecified. A target-
selected tail loss constrains a subset but does not require the remainder of the field to be correct.
Scalar speed removes direction, power removes phase, radial shelling removes angular structure, and
a slope further reduces a spectrum to an exponent. Continuity discrepancy measures a separate
discrete relation involving density evolution, fluxes and boundary treatment. Agreement at one
level need not imply agreement at another.
The smoke is a direct empirical example: peak and gate agreement coexist with poor mean
retention, near-unit field relative error and a maximum absolute velocity divergence exceeding the
target maximum. The CSTR records provide a complementary example: a slope can approach the
8

reference while continuity and energy discrepancies remain large. Reporting these together changes
the scientific interpretation from successful reconstruction to selective metric agreement.
The spectral code change cannot be reduced to an isolated change from vector velocity to speed.
Temporal selection, FFT normalization, shell assignment, fitting range and numerical regulariza-
tion also differ. Variant weights, continuation histories, gradient restrictions and thermodynamic
objectives vary. The association between alignment and improved audit slope is consistent with
the identified contract mismatch, but these records are not a matched causal ablation. Likewise,
the wavelet results do not establish a universal failure of wavelets: they show that the captured
auxiliary detail losses did not establish joint success under the recorded contracts.
Population limits are equally important. The current CSTR training/evaluation overlap pre-
vents treating an unchanged OOD label as a quarantined holdout after adaptation. Consistency
of current reference-slope aggregates is weaker than authenticated historical replay. Row counts
do not establish independent physical trajectories, and repeated candidate panels do not provide
independent replicates. We therefore report neither a pooled pass rate nor a confidence interval or
independent generalization claim.
The original field arrays and checkpoints were not available for this reanalysis. Consequently,
spatial error patterns, casewise replay, unit calibration and exact historical runtime/source bind-
ing remain unresolved. Global-sum agreement, residual discrepancy and proxy tests should not
be relabeled as validated physical conservation, exact equivariance or atomistic topology. A sepa-
rate periodic fluid/magnetic harmonic decoder and a separate temporal generator are outside this
experiment; their results cannot close these gaps.

\section*{6 Conclusion}

In this retrospective scientific-field codec study, recovered extrema and improved fitted spectra did
not establish accurate reconstructed fields. The selected hypersonic smoke fitted its maximum while
retaining little mean kinetic-energy density and exhibiting adverse field and divergence diagnostics.
The distinct full-training checkpoint failed its own spectral assessment despite completion. Audit-
aligned CSTR variants improved the measured slope, but continuity-residual and energy contracts
remained unsatisfied; wavelet-tail variants displayed a different tradeoff rather than joint success.
The practical implication is to publish metric contracts alongside outcomes: observable, spa-
tial and temporal support, transform normalization, shelling, fit range, denominator, reduction,
conditioning information and population relationship. Such specificity makes selective successes in-
terpretable without turning them into unsupported claims of field accuracy or physical validation.

\section*{7 Data and Code A vailability}

Supplementary record S1 is a materialized private export containing captured original source text,
saved reports and manifest metadata, together with exact local checking inputs, executable check-
ing scripts and their saved results. It also contains the primary reference PDFs and inspection
records. The original application source is retained as text; the executable local scripts perform
bounded arithmetic, static metadata and finite-array checks rather than rerunning original training
or decoding.
9

No documented external public or reviewer access mechanism, public archive link, or redistribu-
tion permission is established for this private export. Original field arrays and model checkpoints
are not supplied as replayable artifacts here. The export therefore preserves the evidence used
for the retrospective comparisons but does not establish end-to-end reproducibility of the histor-
ical neural experiments. Artifact identities, variant crosswalks and specific unresolved checks are
retained in the structured verification record.

\section*{References}

[1] V. Saragadam, D. LeJeune, J. Tan, G. Balakrishnan, A. Veeraraghavan, and
R. G. Baraniuk. WIRE: Wavelet Implicit Neural Representations. Proceedings of
the IEEE/CVF Conference on Computer Vision and Pattern Recognition, pp.18507–
18516 (2023). https://openaccess.thecvf.com/content/CVPR2023/html/Saragadam\_WIRE\_
Wavelet\_Implicit\_Neural\_Representations\_CVPR\_2023\_paper.html.
[2] L. Jiang, B. Dai, W. Wu, and C. C. Loy. Focal Frequency Loss for Image Reconstruction
and Synthesis. Proceedings of the IEEE/CVF International Conference on Computer Vision,
pp.13919–13929 (2021). https://openaccess.thecvf.com/content/ICCV2021/html/Jiang\_
Focal\_Frequency\_Loss\_for\_Image\_Reconstruction\_and\_Synthesis\_ICCV\_2021\_paper.
html.
10

\end{document}
